\chapter{ただ一つの速い出力}
% full version: chapters/13_muscles.tex

あなたが世界に向けて速く行うことはすべて、つまり言葉も、視線も、一歩も、筋肉の収縮である。マシンには2つ目の遅い出力、化学的な出力もあり、それについては第16章で扱う。ここでは筋肉の話をする。

連鎖の最後の輪は\nt{ACET}ニューロンである。筋線維の1本1本はちょうど1個のこうしたニューロンに耳を傾け、1個のニューロンは数百から2000本ほどの線維の群れを制御する。面白いことに、神経と筋肉の接続は体の中で最も信頼性が高い。脳の中では接続部が信号の半分以上を失うが、ここではカチッ1回ごとに必要量の数倍の物質が放出される。出力での誤りは動作そのものを損なうので、信頼性は余裕をもって買われているのだ。

\section*{ドラムロールが押す力になる}

カチッ1回で、筋肉は1回短くぴくりと動く。カチッがまばらなら、筋肉は個々のぴくつきとして動く。頻繁なら、ぴくつきは融合してなめらかな力になる。ドラムロールがうなりに溶け合うように。筋肉は、最も単純なデジタル・アナログ変換器のように、スパイク列を力に変える。

\pencilfig{113}{%
% twitch = alpha function; the sum of twitches for spikes 1 s and 0.16 s apart (arbitrary units of time)
\begin{scope}[x=0.95cm, y=0.36cm]
\foreach \dx/\t in {0/まばら,4.6/頻繁}{
  \begin{scope}[xshift=\dx cm]
  \draw[pen] (0,0) -- (4,0); \node[font=\small\itshape, text=pcGray, anchor=south] at (2,5.4) {\t};
  \end{scope}}
\draw[penlite=pcRed] plot[smooth] coordinates {(0.00,0.00) (0.05,0.00) (0.10,0.00) (0.15,0.00) (0.20,0.00) (0.25,0.45) (0.30,0.73) (0.35,0.90) (0.40,0.98) (0.45,1.00) (0.50,0.98) (0.55,0.94) (0.60,0.88) (0.65,0.81) (0.70,0.74) (0.75,0.66) (0.80,0.59) (0.85,0.52) (0.90,0.46) (0.95,0.41) (1.00,0.35) (1.05,0.31) (1.10,0.27) (1.15,0.23) (1.20,0.20) (1.25,0.62) (1.30,0.88) (1.35,1.02) (1.40,1.08) (1.45,1.09) (1.50,1.06) (1.55,1.00) (1.60,0.93) (1.65,0.86) (1.70,0.78) (1.75,0.70) (1.80,0.62) (1.85,0.55) (1.90,0.48) (1.95,0.42) (2.00,0.37) (2.05,0.32) (2.10,0.28) (2.15,0.24) (2.20,0.21) (2.25,0.62) (2.30,0.88) (2.35,1.03) (2.40,1.09) (2.45,1.09) (2.50,1.06) (2.55,1.01) (2.60,0.94) (2.65,0.86) (2.70,0.78) (2.75,0.70) (2.80,0.62) (2.85,0.55) (2.90,0.48) (2.95,0.42) (3.00,0.37) (3.05,0.32) (3.10,0.28) (3.15,0.24) (3.20,0.21) (3.25,0.62) (3.30,0.88) (3.35,1.03) (3.40,1.09) (3.45,1.09) (3.50,1.06) (3.55,1.01) (3.60,0.94) (3.65,0.86) (3.70,0.78) (3.75,0.70) (3.80,0.62) (3.85,0.55) (3.90,0.48) (3.95,0.42) (4.00,0.37)};
\foreach \s in {0.20,1.20,2.20,3.20}{\draw[pcBlue, line width=0.8pt] (\s,-0.35) -- (\s,-1.2);}
\begin{scope}[xshift=4.6cm]
\draw[penlite=pcRed] plot[smooth] coordinates {(0.00,0.00) (0.05,0.00) (0.10,0.00) (0.15,0.00) (0.20,0.00) (0.25,0.45) (0.30,0.73) (0.35,0.90) (0.40,1.35) (0.45,1.68) (0.50,1.85) (0.55,2.19) (0.60,2.51) (0.65,2.64) (0.70,2.84) (0.75,3.13) (0.80,3.22) (0.85,3.27) (0.90,3.55) (0.95,3.62) (1.00,3.54) (1.05,3.82) (1.10,3.88) (1.15,3.80) (1.20,3.98) (1.25,4.06) (1.30,3.97) (1.35,4.07) (1.40,4.16) (1.45,4.09) (1.50,4.10) (1.55,4.23) (1.60,4.17) (1.65,4.10) (1.70,4.26) (1.75,4.23) (1.80,4.07) (1.85,4.27) (1.90,4.27) (1.95,4.12) (2.00,4.26) (2.05,4.29) (2.10,4.17) (2.15,4.24) (2.20,4.31) (2.25,4.21) (2.30,4.21) (2.35,4.31) (2.40,4.24) (2.45,4.16) (2.50,4.31) (2.55,4.27) (2.60,4.10) (2.65,4.30) (2.70,4.29) (2.75,4.15) (2.80,4.28) (2.85,4.31) (2.90,4.18) (2.95,4.25) (3.00,4.32) (3.05,4.22) (3.10,4.21) (3.15,4.32) (3.20,4.25) (3.25,4.17) (3.30,4.31) (3.35,4.27) (3.40,4.11) (3.45,3.86) (3.50,3.56) (3.55,3.25) (3.60,2.93) (3.65,2.63) (3.70,2.33) (3.75,2.06) (3.80,1.81) (3.85,1.58) (3.90,1.38) (3.95,1.19) (4.00,1.03)};
\foreach \s in {0.20,0.36,0.52,0.68,0.84,1.00,1.16,1.32,1.48,1.64,1.80,1.96,2.12,2.28,2.44,2.60,2.76,2.92,3.08,3.24}{\draw[pcBlue, line width=0.8pt] (\s,-0.35) -- (\s,-1.2);}
\end{scope}
\node[lbl, text=pcRed, anchor=west] at (1.2,1.9) {ぴくつき};
\node[lbl, text=pcRed, anchor=south] at (6.6,4.55) {なめらかな力};
\node[lbl, text=pcBlue, anchor=north] at (4.3,-1.35) {ニューロンのカチッ};
\end{scope}
}{まばらなカチッは筋肉の個々のぴくつきを生み、頻繁なカチッは融合して、なめらかではるかに大きな力になる。}

\section*{浮動小数点の力}

脳には力のつまみが2つある。ニューロンがどれだけ頻繁に鳴るかと、線維の群れをいくつ動員するかだ。群れは順番に、最も弱いものから最も強いものへと加わっていく。この順番を計算する者は誰もいない。小さなニューロンほど単に興奮しやすいのだ。その結果、新しく加わる群れはどれも、すでにある力にほぼ同じ\emph{割合}を上乗せする。卵をつまむときは力をごく小刻みに、スーツケースを持ち上げるときは大きな単位で配分する。エンジニアならここに浮動小数点数を見るだろう。桁は動員した群れの数で、細かい値はカチッの頻度で決まる。裏返せば、動きが強いほど精度は下がる。ある有力な仮説によれば、だから私たちの動きはなめらかなのだ。なめらかな指令が最もばらつきを小さくする。

\section*{引けるが押せない}

筋肉は引くことしかできない。だから各関節は、反対向きに引く一対の筋肉が受けもつ。またしても、符号の代わりに2本の配線だ。力の差が関節を回し、力の和が関節を硬くする。なみなみと注いだカップを運ぶとき、あなたは両方の筋肉を同時に緊張させる。回転は同じでも、腕は衝撃にしっかり耐える。

\pencilfig{13}{\pencilart{13}}{筋肉は引くことしかできないので、2つある。力の差が関節を回し、力の和が関節を硬くする。}

\section*{遅れのあるフィードバック}

筋肉には伸張センサーがあり、最も短いループは脊髄の中で直接閉じている。筋肉が引き伸ばされると、その筋肉は自分で収縮し、拮抗筋はその間オフになる。オフにするのは\nt{GLYC}型の抑制性ニューロンだ。この型の仕事がここで見つかった。

だが、どのループにも遅延がある。脊髄では数十ミリ秒、目を経由すると数百ミリ秒。そして遅れたフィードバックは、遅れて道路を見ながらハンドルを切る運転手のように、システムを揺さぶる。ループが長いほど、修正はゆっくり行わなければならない。脊髄ループが振動し始める境目は、毎秒約8回の振動である。これは健康な人なら誰にでもある細かな震えの周波数に近く、伸張ループはその有力な源の一つである。

では、私たちはどうやって速く正確に動けるのか。広く支持されている仮説によれば、予測するのだ。体の内部モデルを保持し、センサーを待たずに指令の結果を計算する。

\section*{メトロノームなしのリズム}

歩行や呼吸はリズミカルだが、そのためのメトロノームはいらない。互いに抑制し合い、しだいに疲れていく2つのニューロン群は、ひとりでに交代で働く。勝者が疲れ、休んだライバルが優位に立ち、それが何度も繰り返される。上からの一定の指令「歩け」が、その場で「左、右」のリズムに変わるのだ。

\fullversion{13}
